\documentclass[10pt,a4paper]{amsart}
\usepackage[margin=19mm]{geometry}
\usepackage{amsmath,amssymb,mathtools}
\usepackage[scr=rsfs,cal=euler]{mathalfa}
\usepackage{xfrac}
\usepackage[unicode,hidelinks,bookmarks=false]{hyperref}
\newcommand{\ie}{i.e.\ }
\newcommand{\cf}{cf.\ }
\newcommand{\resp}{respectively\ }
\newcommand{\Subsection}{\subsection}
\providecommand{\nobreakdash}{\nobreak\hbox{-}}
\setlength{\emergencystretch}{2em}
\multlinegap0pt
\usepackage{amssymb}
\usepackage{bm}
\usepackage{mathtools}
\usepackage{enumerate}
\usepackage{xcolor}
\usepackage{tikz}
\usepackage{float}
\newtheorem{lemma}{Lemma}
\newtheorem{claim}{Claim}
\newcommand{\dif}{\mathop{}\!\mathrm{d}}
\title{\textbf{Maximum spectral sum of graphs}}
\author{Hitesh Kumar, Lele Liu, Hermie Monterde, Shivaramakrishna Pragada, Michael Tait}
\date{Source: arXiv:2604.00512v2 (CC BY 4.0)}
\begin{document}
\maketitle

\noindent\emph{Source and licence.} \textbf{Maximum spectral sum of graphs}. Version arXiv:2604.00512v2 (CC BY 4.0), distributed by arXiv under the Creative Commons Attribution 4.0 International licence, http://creativecommons.org/licenses/by/4.0/, as recorded in the arXiv licence metadata for this exact version and shown on the abstract page https://arxiv.org/abs/2604.00512v2. Full licence notice: \texttt{licenses/arxiv-2604.00512-CC-BY-4.0.txt}.\par\smallskip

\noindent\textit{Editorial scope.} Counted unit: the article's lemma lemma:mixed-cloning (source0.tex lines 605--621). The article's complete original proof of it is delivered with it (source0.tex lines 623--772, one \texttt{proof} environment of 150 lines). No mathematics is written, repaired or re-derived: the statement and the proof are the article's own lines, in the article's own notation, and no retained line is substituted or re-flowed.  Retained uncounted as the source-local context the proof needs: lines 326--331; lines 462--467.  Nothing was omitted from any retained span, so the package carries no omission record.  No retained line contains a citation, so the wrapper carries no bibliography and no bibliography entry.  No retained line contains a figure environment, an image-inclusion command or drawing code, so no image is delivered and the package declares no asset dependency.  Editorial formatting only, in addition: an AMS \texttt{amsart} preamble that restates the article's own declarations and macros used by the retained lines, namely the environments \texttt{lemma}, \texttt{claim} on separate counters, together with the standard packages those lines need.  The retained environments print in this document as Lemma 1, Claim 1 in order of appearance, whereas the article prints the same items with the numbering of its own full document.  The complete native source of the article as distributed in the arXiv e-print for this exact version is bundled unchanged as \texttt{sources/197593/source0.tex}.
\begin{lemma}\label{lemma:spectral_sum_graphon} 
For $K\in \mathcal{K}/\sim$, we have
\[
\sigma(K) = \mu_1(K) + \mu_2(K) = \max_{\substack{\| f\|_2=\| g\|_2 = 1\\ \langle f,g\rangle = 0}} fKf + gKg.
\]
\end{lemma}

\begin{lemma}[Ellipse equation]\label{lemma:ellipse}
For a.e. $x\in [0,1]$, we have    
\[
\mu_1 f(x)^2 + \mu_2 g(x)^2 = \mu_1 + \mu_2. 
\]
\end{lemma}

\begin{lemma}\label{lemma:mixed-cloning}
Let $U\subseteq S^+(g)\cup S^0(g)$ be a measurable subset with positive measure. 
Suppose there exist points $x_1,\ldots,x_k\in U\cap E$ and non-negative real numbers 
$u_1, \ldots, u_k$ with $\sum_i u_i = m(U)$ such that the following two moment constraints are satisfied: 
\begin{equation}\label{eq:moment-constraints}
\sum_{i=1}^k u_i f(x_i)^2= \int_U f(x)^2 \dif{x},
\qquad
\sum_{i=1}^k u_i f(x_i)g(x_i)= \int_U f(x)g(x) \dif{x}, 
\end{equation}
Then the linear moments are also preserved; that is,
\[
\sum_{i=1}^k u_i f(x_i)= \int_U f(x) \dif{x},
\qquad
\sum_{i=1}^k u_i g(x_i)= \int_U g(x) \dif{x}.
\] 
By symmetry, a similar assertion holds if $U\subseteq S^-(g)\cup S^0(g)$.
\end{lemma}

\begin{proof} 
Consider a partition of $U$ into $k$ measurable subsets $U_1, \ldots, U_k$ 
such that $m(U_i) = u_i$ for all $1\le i\le k$. Define a new graphon $\widetilde{W}$ with 
respect to the points $x_i$ and the subsets $U_i$ such that
\[
\widetilde{W}(x,y)=
\begin{cases}
W(x,y), & x,y\notin U;\\
W(x_i,y), & x\in U_i,\ y\notin U; \\
1, & x,y\in U.
\end{cases}
\]
Essentially, we are \emph{cloning} the point $x_i$ over $U_i$. Moreover, since $W(x,y)=1$ 
for all $x,y\in U$, the graphons $W$ and $\widetilde{W}$ agree on $U\times U$. 
Define functions $\tilde{f},\tilde{g}$ using a similar cloning operation, i.e.,  
\begin{equation}\label{eq:tilde-def}
\tilde f(x)=
\begin{cases}
f(x_i), & x\in U_i;\\
f(x), & x\notin U;
\end{cases}
\qquad
\tilde g(x)=
\begin{cases}
g(x_i), & x\in U_i;\\
g(x), & x\notin U.
\end{cases}
\end{equation}

We have the following two claims.

\begin{claim}\label{claim:norm_tilde_f_g}
$\|\tilde{f}\|_2 = \|\tilde{g}\|_2 = 1$, and $\langle \tilde f,\tilde g\rangle = 0$. 
\end{claim}

\begin{proof}[Proof of Claim \ref{claim:norm_tilde_f_g}]
Using \eqref{eq:moment-constraints} we get 
\[
\|\tilde f\|_2^2
= \int_{U^c} f(x)^2 \dif{x} + \sum_{i=1}^k u_i f(x_i)^2
= \left(1-\int_U f(x)^2 \dif{x}\right) + \sum_{i=1}^k u_i f(x_i)^2 = 1.
\]
Here, $U^c = [0,1]\backslash U$. Again, by Lemma \ref{lemma:ellipse}, $g(x_i)^2$ 
is a function of $f(x_i)^2$ given by
\[
g(x_i)^2=\frac{\mu_1+\mu_2-\mu_1 f(x_i)^2}{\mu_2}.
\]
Averaging with weights $u_i$ and using the first constraint in \eqref{eq:moment-constraints} gives
\begin{equation}\label{eq:sum_aver_g2_xi}
\sum_{i=1}^k u_i g(x_i)^2
=\frac{(\mu_1+\mu_2)\, m(U) - \mu_1 \sum u_i f(x_i)^2}{\mu_2}.
\end{equation}
On the other hand, integrating the ellipse equation from Lemma \ref{lemma:ellipse} over $U$ yields
\[
\mu_1 \int_U f(x)^2 \dif{x} + \mu_2\int_U g(x)^2 \dif{x} = (\mu_1+\mu_2)\, m(U),
\] 
which is equivalent to
\begin{equation}\label{eq:integration_g2_over_U}
\int_U g(x)^2 \dif{x} = \frac{(\mu_1+\mu_2)\, m(U) - \mu_1\int_U f(x)^2 \dif{x}}{\mu_2}.
\end{equation}
Combining \eqref{eq:sum_aver_g2_xi} and \eqref{eq:integration_g2_over_U} we see 
$\sum_i u_i g(x_i)^2=\int_U g(x)^2 \dif{x}$. Thus
\[
\|\tilde g\|_2^2
= \int_{U^c} g(x)^2 \dif{x} + \sum_i u_i g(x_i)^2
= 1-\int_U g(x)^2 \dif{x} + \sum_{i=1}^k u_i g(x_i)^2
= 1.
\]

Finally, since $\langle f,g\rangle=0$, we have 
\[
\int_{U^c} f(x)g(x) \dif{x} = -\int_U f(x)g(x) \dif{x},
\]
and hence the second constraint in \eqref{eq:moment-constraints} gives us
\[
\langle \tilde f,\tilde g\rangle
= \int_{U^c} f(x)g(x) \dif{x} + \sum_{i=1}^k u_i f(x_i)g(x_i)=0,
\]
finishing the proof of the claim.
\end{proof}

\begin{claim}\label{claim:tilde_fWf}
We have
\begin{equation}\label{eq:mu_1_W_tilde}
\tilde f\ \widetilde{W}\tilde f = \mu_1+\left(\int_U f(x) \dif{x} -\sum_{i=1}^k u_i f(x_i)\right)^2,
\end{equation}
and 
\begin{equation}\label{eq:mu_2_W_tilde}
\tilde g\ \widetilde{W}\tilde g 
=\mu_2+\left(\int_U g(x) \dif x -\sum_{i=1}^k u_i g(x_i)\right)^2.
\end{equation}
\end{claim}

\begin{proof}[Proof of Claim \ref{claim:tilde_fWf}]
We only prove \eqref{eq:mu_1_W_tilde}; the proof of \eqref{eq:mu_2_W_tilde} is similar. 
Expand $\tilde f\ \widetilde{W}\tilde f$ into the following blocks
\begin{align*}
\tilde f\ \widetilde{W}\tilde f
& =\iint_{U^c\times U^c} W(x,y)f(x)f(y)\,\dif{x} \dif{y} \\
& \quad + 2\sum_{i=1}^k \int_{U_i}\int_{U^c} W(x_i,y)\, f(x_i)\, f(y) \dif{x} \dif{y} \\
& \quad + \iint_{U\times U} \tilde f(x)\tilde f(y) \dif{x} \dif{y}.
\end{align*}
Note that the second line equals $2\sum_i u_i f(x_i)\int_{U^c} W(x_i,y)f(y)\,\dif y$, while the third line can be written as
$\big(\int_U \tilde f(x) \dif x \big)^2 = \big(\sum_i u_i f(x_i)\big)^2$ by \eqref{eq:tilde-def}. So
\begin{equation}\label{eq:Qmix-f}
\tilde f\ \widetilde{W}\tilde f
= I+2\sum_{i=1}^k u_i f(x_i)\int_{U^c} W(x_i,y)f(y)\,\dif{y} + \left(\sum_{i=1}^k u_i f(x_i)\right)^2,
\end{equation}
where $I:=\iint_{U^c\times U^c} W(x,y) f(x) f(y) \dif{x} \dif{y}$. Using the 
eigenvalue-equation at $x\in U\cap E$, we have 
\begin{equation}\label{eq:row-outside}
\int_{U^c} W(x,y)f(y)\,\dif{y} = \mu_1 f(x)-\int_U f(y) \dif{y}.
\end{equation}
Using \eqref{eq:row-outside} at $x=x_i$ we rewrite the mixed term in \eqref{eq:Qmix-f}:
\[
2\sum_{i=1}^k u_i f(x_i) \int_{U^c}W(x_i,y)f(y)\,\dif{y} 
= 2\mu_1\sum_{i=1}^k u_i f(x_i)^2 - 2\sum_{i=1}^k u_i f(x_i)\int_U f(y) \dif{y}. 
\]
Combining the above equation with \eqref{eq:Qmix-f} yields
\begin{equation}\label{eq:Qmix-f-2}
\tilde f\, \widetilde{W}\tilde f
= I + 2\mu_1\sum_{i=1}^k u_i f(x_i)^2-2\sum_{i=1}^k u_i f(x_i)\int_U f(x)\dif{x}+\left(\sum_{i=1}^k u_i f(x_i)\right)^2.
\end{equation}

Now, we compute the term $\mu_1=f\, Wf$ using the same block expansion. 
Applying \eqref{eq:row-outside} for $x\in U\cap E$, we obtain
\begin{align*}
\mu_1 
& = I + 2\int_U f(x)\left(\mu_1 f(x)-\int_U f(y)\dif{y}\right) \dif{x} + \left(\int_U f(x)\dif{x}\right)^2 \\
& = I + 2\mu_1\int_U f(x)^2 \dif{x} - \left(\int_U f(x) \dif{x}\right)^2.
\end{align*}
Subtract this identity from \eqref{eq:Qmix-f-2} to get
\[
\tilde f\ \widetilde{W}\tilde f -\mu_1
=2\mu_1\left(\sum_{i=1}^k u_i f(x_i)^2 - \int_U f(x)^2 \dif{x} \right)
+\left(\int_U f(x)\dif{x} -\sum_{i=1}^k u_i f(x_i)\right)^2.
\]
By \eqref{eq:moment-constraints}, the first term vanishes, and we get the desired claim. 
\end{proof}

Because $\tilde f,\tilde g$ are orthonormal, using Lemma \ref{lemma:spectral_sum_graphon}, we have
\begin{align*}
\sigma(\widetilde{W})
& \geq \tilde f\ \widetilde{W}\tilde f + \tilde g\ \widetilde{W}\tilde g \\
& = \mu_1+\mu_2 + \left(\int_U f(x) \dif{x} -\sum_i u_i f(x_i)\right)^2 
+ \left(\int_U g(x) \dif{x} -\sum_i u_i g(x_i)\right)^2.
\end{align*}
Since $W$ is extremal, $\sigma(W)\ge \sigma(\widetilde{W})$, and so the integrals in the 
last inequality must be zero. The assertion follows.
\end{proof}

\end{document}
